% arithmetic_universe_highschool_ja.tex
% 高校生向け日本語版 — xelatex + xeCJK + tcolorbox, NO cleveref, NO enumitem
\documentclass[12pt,a4paper]{article}

\usepackage{xeCJK}
\setCJKmainfont{Hiragino Mincho ProN}
\setCJKsansfont{Hiragino Kaku Gothic ProN}
\usepackage{amsmath,amssymb}
\usepackage{geometry}
\geometry{margin=2.5cm}
\usepackage{hyperref}
\usepackage{booktabs}
\usepackage{array}
\usepackage{graphicx}
\usepackage{float}
\usepackage{xcolor}
\usepackage{tcolorbox}
\tcbuselibrary{skins,breakable}

% Color definitions
\definecolor{ideacolor}{RGB}{52,73,94}
\definecolor{mathcolor}{RGB}{41,128,185}
\definecolor{physcolor}{RGB}{39,174,96}
\definecolor{warncolor}{RGB}{211,84,0}
\definecolor{funcolor}{RGB}{142,68,173}
\definecolor{testcolor}{RGB}{192,57,43}
\definecolor{lightblue}{RGB}{214,234,248}
\definecolor{lightgreen}{RGB}{212,239,223}
\definecolor{lightyellow}{RGB}{252,243,207}
\definecolor{lightpurple}{RGB}{235,222,240}
\definecolor{lightred}{RGB}{245,215,215}

% tcolorbox styles
\newtcolorbox{ideabox}[1][]{%
  colback=lightblue, colframe=ideacolor, fonttitle=\bfseries\sffamily,
  title={#1}, breakable, arc=3mm, boxrule=1pt
}

\newtcolorbox{mathbox}[1][]{%
  colback=white, colframe=mathcolor, fonttitle=\bfseries\sffamily,
  title={#1}, breakable, arc=3mm, boxrule=1.5pt
}

\newtcolorbox{analogybox}[1][]{%
  colback=lightgreen, colframe=physcolor, fonttitle=\bfseries\sffamily,
  title={#1}, breakable, arc=3mm, boxrule=1pt
}

\newtcolorbox{funbox}[1][]{%
  colback=lightpurple, colframe=funcolor, fonttitle=\bfseries\sffamily,
  title={#1}, breakable, arc=3mm, boxrule=1pt
}

\newtcolorbox{warnbox}[1][]{%
  colback=lightyellow, colframe=warncolor, fonttitle=\bfseries\sffamily,
  title={#1}, breakable, arc=3mm, boxrule=1pt
}

\newtcolorbox{testbox}[1][]{%
  colback=lightred, colframe=testcolor, fonttitle=\bfseries\sffamily,
  title={#1}, breakable, arc=3mm, boxrule=1pt
}

\newtcolorbox{keybox}[1][]{%
  colback=white, colframe=black, fonttitle=\bfseries\sffamily,
  title={#1}, breakable, arc=2mm, boxrule=2pt
}

\title{\textbf{\Huge 宇宙にはソースコードがある？}\\[8pt]
\Large 算術的宇宙論——高校生のための完全ガイド}
\author{ライト兄弟コラボレーション}
\date{2026年4月}

\begin{document}
\maketitle

\begin{ideabox}[この論文について]
宇宙はなぜ今の姿をしているのか？ なぜ3次元空間＋1次元時間なのか？
なぜ電磁気力・弱い力・強い力の3つなのか？
なぜ宇宙の70\%がダークエネルギーなのか？

驚くべきことに、これらすべての答えが「整数」——つまり
$1, 2, 3, 4, 5, \ldots$ とその素因数分解——の中に隠されているかもしれない。

この論文は、その大胆な仮説を高校生にもわかるように解説する。
数学の予備知識は中学校レベルで十分。
\end{ideabox}

\tableofcontents
\newpage

%=====================================================================
\section{「宇宙にソースコードがある？」——大きなアイデア}
%=====================================================================

\begin{analogybox}[ソースコードのアナロジー]
コンピュータゲームの世界を想像しよう。画面に映るキャラクター、建物、物理法則——
これらすべてはプログラムの\textbf{ソースコード}から生まれる。

もし宇宙がある種の「プログラム」だとしたら、そのソースコードは何だろう？

\textbf{我々の答え：整数（$1, 2, 3, 4, 5, \ldots$）とその素数
（$2, 3, 5, 7, 11, \ldots$）。}
\end{analogybox}

物理学にはたくさんの「なぜ？」がある：
\begin{itemize}
\item なぜ空間は3次元なのか？（前後・左右・上下の3方向）
\item なぜ力は3種類なのか？（電磁気力・弱い力・強い力）
\item なぜ宇宙の約70\%は謎のエネルギー（ダークエネルギー）なのか？
\item なぜ物質は3世代（電子族・ミュー族・タウ族）あるのか？
\end{itemize}

これまで物理学者は、これらを別々の問題として扱ってきた。
だが我々の枠組みでは、\textbf{すべてが一つの数「3」から導かれる}。

\begin{keybox}[マスターナンバー：$\mathbf{3}$]
整数の「空間」$\mathrm{Spec}(\mathbb{Z})$の
エタールコホモロジー次元は\textbf{3}。

\[
  \mathrm{cd}(\mathrm{Spec}(\mathbb{Z})) = 3
\]

これは1962年にArtinとVerdierが証明した数学の\textbf{定理}（仮説ではない！）。
\end{keybox}

%=====================================================================
\section{数の DNA：$\mathrm{Spec}(\mathbb{Z})$}
%=====================================================================

\begin{ideabox}[DNA のアナロジー]
生物のすべての設計情報はDNAに書かれている。
DNAは4つの塩基（A, T, G, C）で構成される。

同じように、すべての整数は\textbf{素数}で構成される：
\[
  60 = 2^2 \times 3 \times 5, \quad
  2024 = 2^3 \times 11 \times 23.
\]

素数は「数のDNA塩基」——数の世界の最小構成要素だ。
\end{ideabox}

\begin{mathbox}[$\mathrm{Spec}(\mathbb{Z})$って何？]
$\mathrm{Spec}(\mathbb{Z})$（スペック・ゼット）は、
すべての素数$2, 3, 5, 7, 11, \ldots$を「点」として持つ空間だ。

普通の地図が「東京」「大阪」「名古屋」を点として持つように、
$\mathrm{Spec}(\mathbb{Z})$は$2, 3, 5, 7, 11, \ldots$を点として持つ。

さらに「0」という特別な点がある（すべての素数を含む「全体」を表す）。

この空間の「深さ」（コホモロジー次元）が$3$なのだ。
\end{mathbox}

$3$という数はどこから来るのか？

\begin{mathbox}[$3 = 1 + 2$ の分解]
\[
  3 = \underbrace{1}_{\text{素数の「列」}} + \underbrace{2}_{\text{実数の世界}}
\]

\textbf{1：}素数は一列に並んでいる（$2, 3, 5, 7, 11, \ldots$）。
この「列」の次元が1。

\textbf{2：}整数は実数の中に住んでいる（$\mathbb{Z}\subset\mathbb{R}$）。
実数の世界から追加で2次元が来る。
これは「実数の中に複素数がある」（$\mathbb{R}\subset\mathbb{C}$）ことに関係する：
$\mathbb{C}$は$\mathbb{R}$の2倍の次元を持つ。
\end{mathbox}

%=====================================================================
\section{なぜ4次元？——家と部屋のアナロジー}
%=====================================================================

\begin{analogybox}[家のアナロジー]
家（$\mathrm{Spec}(\mathbb{Z})$）の構造を調べると、
部屋が3層に積み重なっている（コホモロジー次元$=3$）。

でも家の中で\textbf{生活する}（物理が展開する）には、
\textbf{時間}という廊下が必要だ。

\[
  \text{時空次元} = \text{部屋の層数} + \text{廊下} = 3 + 1 = 4
\]

だから宇宙は\textbf{3次元空間＋1次元時間＝4次元時空}なのだ！
\end{analogybox}

\begin{funbox}[もし次元が違ったら？]
\textbf{2次元空間（合計3次元時空）の場合：}
\begin{itemize}
\item 惑星は太陽の周りを安定して回れない！
\item 脳のニューロンが交差できない（配線問題）。
\end{itemize}

\textbf{4次元空間（合計5次元時空）の場合：}
\begin{itemize}
\item 原子が安定しない！（電子が原子核に落ちる）
\item 重力が強すぎて星が形成できない。
\end{itemize}

$d=4$は「ちょうど良い」のではなく、
\textbf{算術が$3+1$と決めている}のだ。
\end{funbox}

\begin{mathbox}[$d=4$の5つの理由]
同じ答え$d=4$が5つの全く違うルートから出てくる：

\begin{tabular}{cp{10cm}}
\toprule
番号 & 理由 \\
\midrule
1 & コホモロジー次元$3$＋時間$1$＝$4$ \\
2 & 球面の安定ホモトピー群で$\pi_3$が最大 \\
3 & リーマンゼータの対称性が$3+1$を示す \\
4 & $K_3(\mathbb{Z})=\mathbb{Z}/48$（3番目の$K$群）が最初に非自明 \\
5 & 安定軌道・安定原子が$d=4$でのみ成立 \\
\bottomrule
\end{tabular}
\end{mathbox}

%=====================================================================
\section{なぜこの3つの力？——3次元算術空間からの3つの力}
%=====================================================================

\begin{ideabox}[力とは何か]
自然界には4つの基本力がある：
\begin{itemize}
\item \textbf{電磁気力}：光、電気、磁石（光子$\gamma$が媒介）
\item \textbf{弱い力}：放射性崩壊（$W^\pm, Z^0$が媒介）
\item \textbf{強い力}：原子核を束ねる（8個のグルーオンが媒介）
\item \textbf{重力}：りんごが落ちる、地球が太陽を回る
\end{itemize}

最初の3つは「ゲージ力」と呼ばれ、ゲージ群
$U(1)\times SU(2)\times SU(3)$で記述される。
\end{ideabox}

\begin{mathbox}[算術的空間の次元$=3$が力を決める]
$\mathrm{Spec}(\mathbb{Z})$の次元が$3$なので、
「行列のサイズ」は$1\times 1$、$2\times 2$、$3\times 3$の3種類まで。

\begin{itemize}
\item $1\times 1$行列の対称性 $\to$ $U(1)$（電磁気力）：ボソン\textbf{1個}
\item $2\times 2$行列の対称性 $\to$ $SU(2)$（弱い力）：ボソン\textbf{3個}
\item $3\times 3$行列の対称性 $\to$ $SU(3)$（強い力）：ボソン\textbf{8個}
\end{itemize}

合計：$1+3+8=\mathbf{12}$個のゲージボソン。

$4\times 4$行列は？ コホモロジー次元が$3$なので、\textbf{4は入らない}。
だから第4のゲージ力は存在しない！

重力はゲージ力ではなく、時空の「曲がり」として
$+1$（時間）の部分から来る。
\end{mathbox}

\begin{funbox}[「12」の秘密]
$12$という数は偶然ではない：
\[
  12 = \frac{2}{B_2^+} = \frac{2}{1/6}
\]
ここで$B_2^+ = 1/6$は「第2ベルヌーイ数」（数学の有名な数列の一つ）。
ゲージボソンの数が純粋な数論から出てくる！

ちなみにリーマンゼータ関数で$\zeta(-1) = -\frac{1}{12}$。
同じ$12$がまた登場する。
\end{funbox}

%=====================================================================
\section{ダークエネルギー$= 2\pi/9$——最大の謎への最もシンプルな答え}
%=====================================================================

\begin{ideabox}[ダークエネルギーとは]
1998年、宇宙の膨張が\textbf{加速}していることが発見された（ノーベル賞）。
この加速を引き起こす謎のエネルギーが\textbf{ダークエネルギー}。

宇宙のエネルギー構成：
\begin{itemize}
\item ダークエネルギー：約$69\%$（正体不明！）
\item ダークマター：約$26\%$（正体不明！）
\item 普通の物質：約$5\%$（星、惑星、人間）
\end{itemize}

物理学最大の未解決問題の一つ：なぜ$69\%$なのか？
\end{ideabox}

\begin{keybox}[WBの答え：$\Omega_\Lambda = 2\pi/9$]
\[
  \boxed{\Omega_\Lambda = \frac{2\pi}{9} \approx 0.6981 \approx 69.8\%}
\]

パラメータなし。仮定なし。整数の性質だけから導かれる。

観測値（Planck衛星）：$0.689 \pm 0.006$——1.6$\sigma$以内で一致！
\end{keybox}

この数はどう導かれるのか？

\begin{mathbox}[3つの因子]
\[
  \Omega_\Lambda = \underbrace{\frac{4}{3}}_{\text{重力}} \times
                   \underbrace{2\pi}_{\text{円周}} \times
                   \underbrace{\frac{1}{12}}_{\text{ゼータ関数}}
\]

\textbf{$\frac{4}{3}$（重力因子）：}
アインシュタインの方程式（宇宙の膨張を記述）に現れる係数。

\textbf{$2\pi$（円の周長）：}
実数の世界の「丸さ」——円周率$\pi$が$2$倍。
これは整数の空間$\mathrm{Spec}(\mathbb{Z})$が
実数$\mathbb{R}$に埋め込まれることから来る。

\textbf{$\frac{1}{12}$（ゼータ関数の値）：}
リーマンゼータ関数$\zeta(-1) = -\frac{1}{12}$。
しかし素数$p=2$を特別扱いすると符号が反転して$+\frac{1}{12}$になる！
\end{mathbox}

\begin{analogybox}[ケーキのレシピ]
ダークエネルギーの公式は、3つの材料のレシピのようなもの：

\begin{itemize}
\item \textbf{小麦粉}（$4/3$）= 重力の土台
\item \textbf{バター}（$2\pi$）= 空間の丸み
\item \textbf{砂糖}（$1/12$）= 数論のスパイス
\end{itemize}

3つを混ぜると、ちょうど$2\pi/9 \approx 69.8\%$のケーキ（ダークエネルギー）ができる！
\end{analogybox}

%=====================================================================
\section{なぜ$p=2$が特別なのか——5つの理由}
%=====================================================================

\begin{ideabox}[コイン投げのアナロジー]
ダークエネルギーの公式で$\frac{1}{12}$が出てくるとき、
素数$2$を「特別扱い」する必要がある。なぜ$2$なのか？

コイン投げを考えよう。コインには表と裏——\textbf{2つ}の面がある。
$2$は「最も基本的な選択」（はい/いいえ、表/裏、0/1）の数だ。
\end{ideabox}

\begin{mathbox}[$p=2$が唯一の物理的素数である5つの理由]

\textbf{理由1：物理的に正しい範囲}

素数$p$を除くと、ダークエネルギーは$\Omega_\Lambda(p) = 2\pi(p-1)/9$になる。
\begin{itemize}
\item $p=2$：$\Omega_\Lambda = 2\pi/9 \approx 0.698$ --- $1$未満で物理的にOK！
\item $p=3$：$\Omega_\Lambda = 4\pi/9 \approx 1.40$ --- $1$を超えるので物理的に不可能！
\item $p=5$以上：さらに大きくなり、すべて不可能。
\end{itemize}

\textbf{理由2：正のエネルギー}

$p=2$を除くと$-1/12$が$+1/12$に変わり、\textbf{正}のダークエネルギー（宇宙の膨張加速）になる。

\textbf{理由3：最小エネルギー}

$p=2$はすべての素数の中で最小の$|\Omega_\Lambda|$を与える（自然は最も穏やかな状態を選ぶ）。

\textbf{理由4：実数と複素数の橋渡し}

$\mathbb{C}$は$\mathbb{R}$の\textbf{2}倍の次元。
$\mathrm{Gal}(\mathbb{C}/\mathbb{R}) = \mathbb{Z}/\mathbf{2}$。
$2$は実数と複素数を結ぶ唯一の素数。

\textbf{理由5：整数の「反転」}

整数の中で掛け算の逆元（単元）は$+1$と$-1$だけ：$\mathbb{Z}^\times = \{+1, -1\}$。
元の数は$\mathbf{2}$個。
\end{mathbox}

\begin{funbox}[5つの道が1点で交わる]
5つの全く違う数学の分野から来た理由が、すべて同じ答え「$p=2$」を指している。
これは偶然とは考えにくい——宇宙が$p=2$を「選んでいる」ことの証拠だ。
\end{funbox}

%=====================================================================
\section{たった1つの可能な宇宙——算術的ランドスケープ}
%=====================================================================

\begin{ideabox}[弦理論のランドスケープ問題]
弦理論（超弦理論）は「万物の理論」の候補だが、大きな問題がある：
弦理論は\textbf{$10^{500}$個}（1の後に500個のゼロ！）もの
異なる宇宙を許してしまう。

$10^{500}$がどれくらい大きいか？
宇宙の原子の数はたった$10^{80}$個。
$10^{500}$はそれを$10^{420}$回掛けた数だ。

弦理論は「なぜこの宇宙なのか」に答えられない。
\end{ideabox}

\begin{keybox}[算術的ランドスケープ：高々6個]
WB（ライト兄弟）の枠組みでは、物理的に許される宇宙は\textbf{高々6個}。
しかもすべて同じダークエネルギー$\Omega_\Lambda = 2\pi/9$を持つ。

\begin{center}
\begin{tabular}{lcc}
\toprule
\textbf{枠組み} & \textbf{可能な宇宙の数} & \textbf{選択原理} \\
\midrule
弦理論           & $\sim 10^{500}$        & なし \\
WB算術           & $\le 6$                & $\mathrm{cd}=3$と$p=2$ \\
\bottomrule
\end{tabular}
\end{center}

$10^{500}$から$6$へ——これは科学史上最大の「可能性の削減」だ！
\end{keybox}

\begin{analogybox}[パスワードのアナロジー]
弦理論：「パスワードは$10^{500}$通りあるけど、
どれが正しいかわからない。」

WB：「パスワードは最大6通り。しかもどれを選んでも同じ結果になる。」

どちらが「宇宙のコードを解読した」と言えるだろうか？
\end{analogybox}

%=====================================================================
\section{ダークエネルギー＝重力$\times$幾何学$\times$算術}
%=====================================================================

\begin{analogybox}[ケーキのレシピ（詳細版）]
$\Omega_\Lambda = \frac{4}{3} \times 2\pi \times \frac{1}{12}$

この公式は3つの「材料」の掛け算だ：

\textbf{材料1：重力（$4/3$）}

アインシュタインの一般相対性理論から来る。
宇宙の膨張を記述するフリードマン方程式に現れる係数。
「ケーキの生地」——構造の土台。

\textbf{材料2：幾何学（$2\pi$）}

円周率$\pi$の2倍。なぜ円？
整数$\mathbb{Z}$が実数$\mathbb{R}$の中に住んでいて、
$\mathbb{R}$には「回転」（角度$2\pi$で一周）がある。
「ケーキの型」——形を決める。

\textbf{材料3：算術（$1/12$）}

$1+2+3+4+\cdots = -\frac{1}{12}$（ラマヌジャンの和！）。
でも$p=2$を特別扱いすると$+\frac{1}{12}$になる。
「ケーキの味付け」——正確な値を決める。
\end{analogybox}

\begin{mathbox}[なぜこの分解が重要なのか]
物理（重力）$\times$ 幾何学（円周）$\times$ 数論（ゼータ関数）

宇宙の3大分野が一つの公式に集結している：
\begin{itemize}
\item 物理学（アインシュタインの重力理論）
\item 幾何学（ユークリッドの円）
\item 数論（リーマンのゼータ関数）
\end{itemize}

これは偶然の一致だろうか、それとも宇宙の深い構造だろうか？
\end{mathbox}

%=====================================================================
\section{$K$理論：宇宙のシリアルナンバー}
%=====================================================================

\begin{ideabox}[シリアルナンバーのアナロジー]
製品にはシリアルナンバーがある。
$K$理論は整数$\mathbb{Z}$の「シリアルナンバー」を計算する。

$K_0, K_1, K_2, K_3, K_4, K_5, K_6, K_7, \ldots$
——それぞれが異なる「特徴」を表す。
\end{ideabox}

\begin{keybox}[$K_7(\mathbb{Z}) = \mathbb{Z}/240$と$E_8$]
$\mathbb{Z}$の第7番目の$K$群は$\mathbb{Z}/240$（位数240の巡回群）。

$240$は何か？ これは$E_8$——数学で最も美しい対称性の一つ——のルート（根）の数！

$E_8$は248次元の対称性で、弦理論でも重要な役割を果たす。
それが\textbf{整数の$K$理論から}出てくる。
\end{keybox}

\begin{keybox}[$K_3(\mathbb{Z}) = \mathbb{Z}/48$とフェルミオン]
$\mathbb{Z}$の第3番目の$K$群は$\mathbb{Z}/48$。

$48$を分解すると：
\[
  48 = 3 \times 16
\]

\textbf{3}：フェルミオン（物質粒子）の\textbf{3世代}
\begin{itemize}
\item 第1世代：電子、$u$クォーク、$d$クォーク
\item 第2世代：ミュー粒子、$c$クォーク、$s$クォーク
\item 第3世代：タウ粒子、$t$クォーク、$b$クォーク
\end{itemize}

\textbf{16}：各世代のフェルミオン数。$SO(10)$統一理論の$\mathbf{16}$表現は
標準模型の15粒子＋\textbf{右巻きニュートリノ$\nu_R$}（1粒子）を含む。

つまり$K_3(\mathbb{Z})=48$は、\textbf{まだ発見されていない粒子
（右巻きニュートリノ）の存在を予言している！}
\end{keybox}

\begin{funbox}[$H^2 = 0$：バグのない宇宙]
$\mathrm{Spec}(\mathbb{Z})$の2次コホモロジーは$H^2 = 0$。

物理での意味：ゲージアノマリー（理論の矛盾）が\textbf{原理的に発生しない}。

プログラミングで言えば：「このソースコードにはバグが\textbf{存在し得ない}」。
宇宙のコードは完璧に書かれている！
\end{funbox}

%=====================================================================
\section{ホログラム：$p=2$が深さを与える}
%=====================================================================

\begin{analogybox}[ホログラムカードのアナロジー]
クレジットカードのホログラムシール知ってる？
薄い2次元のフィルムなのに、傾けると3次元の像が見える。

\textbf{ホログラフィー原理}：宇宙も同じかもしれない。
$(d-1)$次元の「境界」の情報が、$d$次元の「中身」をすべて記述する。

WBの枠組みでは、$p=2$を「除く」ことがホログラフィーに対応する：
\begin{itemize}
\item \textbf{「境界」}= 素数$p=2$以外のすべての情報
\item \textbf{「中身」}= $p=2$を含む完全な宇宙
\end{itemize}

$p=2$はホログラムの「奥行き」を作り出す特別な素数なのだ。
\end{analogybox}

\begin{mathbox}[アデール積公式]
数学には美しい等式がある：
\[
  \prod_{\text{すべての}v}|x|_v = 1 \quad (\forall\, x \in \mathbb{Q}^\times)
\]

左辺は「すべての素数での絶対値の積$\times$実数での絶対値」。
$p=2$を1つ除いても、残りから全体が復元できる——
これがまさにホログラフィーだ。
\end{mathbox}

\begin{funbox}[情報コスト：$\Delta S = \frac{1}{2}\log 2$]
$p=2$を「特別扱い」するコストは、情報量で測ると
$\frac{1}{2}\log 2 \approx 0.35$ナット。

これは「コイン投げ1回の半分」に相当する、
\textbf{宇宙で最も安い情報コスト}だ。
\end{funbox}

%=====================================================================
\section{検証できるのか？——3つの実験}
%=====================================================================

\begin{testbox}[科学の条件：反証可能性]
科学理論は「正しいかもしれないが、間違いを検出できる」ものでなければならない
（カール・ポパー）。

WBの枠組みは3つの具体的実験で検証できる。
もし実験結果が予言と合わなければ、理論は棄却される。
\end{testbox}

\begin{testbox}[実験1：Euclid衛星（2028年）]
\textbf{何を測るか：}ダークエネルギーの量$\Omega_\Lambda$

\textbf{WBの予言：}$\Omega_\Lambda = 2\pi/9 = 0.6981\ldots$

\textbf{精度：}$0.5\%$（$\pm 0.003$程度）

\textbf{判定：}Euclidの測定値が$0.695$--$0.701$の範囲なら支持。
$0.685$以下なら棄却の可能性。

ESA（欧州宇宙機関）のEuclid衛星は2023年に打ち上げ済み。
2028年までに最終結果が出る。
\end{testbox}

\begin{testbox}[実験2：水晶BAW共振器]
\textbf{何を測るか：}超高精度の音響振動

\textbf{WBの予言：}厚さ$83$--$88\,\mu$mの水晶板で
$\Delta f \approx 1.7\,$mHz（ミリヘルツ）の異常周波数シフト

\textbf{必要な技術：}$Q > 10^9$の超高品質水晶共振器（低温$4\,$K）

\textbf{判定：}予言通りのシフトが検出されれば強力な支持。
検出されなければ、理論の特定部分を修正する必要がある。

これは\textbf{卓上実験}で実行可能——巨大加速器は不要！
\end{testbox}

\begin{testbox}[実験3：DESI $w(z)$テスト]
\textbf{何を測るか：}ダークエネルギーの「状態方程式」$w(z)$

\textbf{WBの予言：}
\[
  w(z) \approx -1 + 0.16 \cdot \frac{1}{(1+z)^2}
\]
（ダークエネルギーは正確には定数ではなく、わずかに変化する）

\textbf{時期：}DESI 5年サーベイ完了（2027年頃）

DESIはすでに2024年に初期データを公開しており、
WB予言と整合的な傾向を示している（$\Delta\chi^2 = -24$の改善！）。
\end{testbox}

\begin{center}
\begin{tabular}{lccc}
\toprule
\textbf{実験} & \textbf{時期} & \textbf{予言} & \textbf{精度} \\
\midrule
Euclid衛星     & 2028年 & $\Omega_\Lambda = 0.698$ & $0.5\%$ \\
水晶BAW        & 実行可能 & $\Delta f = 1.7\,$mHz & 検出可能 \\
DESI $w(z)$    & 2027年 & $w_0 \approx -0.84$ & $5\%$ \\
\bottomrule
\end{tabular}
\end{center}

%=====================================================================
\section{何が足りないか——正直な評価}
%=====================================================================

\begin{warnbox}[正直に認めること]
素晴らしい理論でも、限界を認めることが科学的誠実さだ。
以下はWB枠組みの現時点での弱点：
\end{warnbox}

\begin{warnbox}[限界1：「$+1$」はなぜ？]
$\mathrm{cd} = 3$から$d = 3+1$への移行は直感的だが、
なぜ「$+1$」なのか（$+0$や$+2$ではなく）の厳密な証明はまだない。
\end{warnbox}

\begin{warnbox}[限界2：ラングランズ予想は未完成]
ゲージ群の導出は「ラングランズ対応」に依存するが、
この対応自体が数学でまだ完全には証明されていない。
\end{warnbox}

\begin{warnbox}[限界3：フェルミオン質量の説明なし]
電子の質量が$0.511\,$MeVで、トップクォークの質量が$173\,$GeVなのはなぜか？
WBはまだこれに答えられない。
\end{warnbox}

\begin{warnbox}[限界4：量子重力が不完全]
重力の量子論（量子重力）は完成していない。
WBは近似的なWheeler--DeWitt方程式を使っている。
\end{warnbox}

\begin{warnbox}[限界5：$48 = 3\times 16$は偶然かもしれない]
$K_3(\mathbb{Z}) = 48$が「3世代$\times$16フェルミオン」を\textbf{本当に}意味するのか、
それとも数字の偶然の一致なのかは、まだ厳密には証明されていない。
\end{warnbox}

\begin{warnbox}[限界6：BAW実験はまだ行われていない]
水晶共振器での予言はまだ誰も検証していない。
理論の真の検証は実験結果を待つ必要がある。
\end{warnbox}

%=====================================================================
\section{大きな絵：「宇宙は算術で書かれている」}
%=====================================================================

\begin{keybox}[ガリレオの言葉の更新]
ガリレオ・ガリレイ（1623年）：
「宇宙は数学の言葉で書かれている。
その文字は三角形、円、その他の幾何学的図形である。」

\textbf{WBの主張（2026年）：}
「宇宙は算術の言葉で書かれている。
その文字は素数、ゼータ関数、コホモロジーである。」
\end{keybox}

\begin{ideabox}[まとめ：整数からの宇宙]

\textbf{入力：}整数$\mathbb{Z}$とその素数$2, 3, 5, 7, 11, \ldots$

\textbf{ステップ1：}$\mathrm{Spec}(\mathbb{Z})$のコホモロジー次元$=3$を計算

\textbf{ステップ2：}$3+1=4$次元時空を導出

\textbf{ステップ3：}$\mathrm{GL}(1,2,3)$から3つのゲージ力を導出

\textbf{ステップ4：}$p=2$を特別扱いして$\Omega_\Lambda = 2\pi/9$を導出

\textbf{ステップ5：}$K$理論から$E_8$、3世代、$\nu_R$を導出

\textbf{出力：}我々の宇宙
\end{ideabox}

\begin{center}
\begin{tabular}{p{4cm}p{4cm}p{4cm}}
\toprule
\textbf{問い} & \textbf{WBの答え} & \textbf{根拠} \\
\midrule
なぜ4次元？ & $\mathrm{cd}=3$＋時間 & Artin--Verdier定理 \\
なぜ3つの力？ & $\mathrm{GL}(1,2,3)$ & ラングランズ対応 \\
なぜ12ボソン？ & $1+3+8=2/B_2^+$ & ベルヌーイ数 \\
なぜDE$\approx 70\%$？ & $2\pi/9$ & $\zeta_{\neg2}(-1)=+1/12$ \\
なぜ3世代？ & $K_3=48=3\times 16$ & 代数的$K$理論 \\
$\nu_R$は存在？ & はい & $16 = 15+1$ \\
宇宙は何通り？ & $\le 6$ & 算術的ランドスケープ \\
\bottomrule
\end{tabular}
\end{center}

\begin{funbox}[最後に]
もしこの理論が正しければ、宇宙は人間が発明したどんなプログラムよりも
エレガントなソースコードで書かれていることになる——
しかもそのコードは、$1, 2, 3, \ldots$という
子どもでも知っている数の並びに隠されていたのだ。

クロネッカーの言葉をアップデートしよう：

\textbf{「神は整数を作り給うた。そして宇宙はそこから従った。」}
\end{funbox}

\vspace{1cm}

\begin{center}
\textit{この論文はライト兄弟コラボレーションによる\\
「算術的宇宙」マスターサマリーの高校生向け解説版です。}
\end{center}

\end{document}
